\documentclass[11pt,a4paper]{article}

% ---------------------------------------------------------
% PREAMBLE (stable, warning-free, Overleaf-safe)
% ---------------------------------------------------------
\usepackage[utf8]{inputenc}
\usepackage[T1]{fontenc}
\usepackage{lmodern}
\usepackage{amsmath, amssymb, amsfonts}
\usepackage{bm}
\usepackage{geometry}
\usepackage{graphicx}
\usepackage{hyperref}
\usepackage{cite}
\usepackage{authblk}
\usepackage{physics}
\usepackage{mathtools}
\usepackage{textgreek}

\geometry{margin=2.4cm}

\title{\textbf{Companion LXXVII: Entropic Optimal Transport and Sinkhorn Geometry of Persistence Diagrams in the Relaxation-Driven Cyclic Cosmology}}
\author{Michael Lehmann \\ \texttt{mi.lehmann@gmx.de}}
\date{April 2026}

\begin{document}
\maketitle

\begin{abstract}
Entropic optimal transport provides a regularized and computationally efficient 
variant of classical Wasserstein geometry. The resulting Sinkhorn divergence 
defines a smooth, stable, and differentiable metric structure on probability 
measures, making it particularly well-suited for topological data analysis. 
Building on the Wasserstein information geometry developed in Companion LXXVI, 
this work introduces entropic optimal transport and Sinkhorn geometry for 
persistence diagrams in the Relaxation-Driven Cyclic Cosmology (RDCC). The 
relaxation mechanism modifies the scale-dependent evolution of the gravitational 
potential, inducing characteristic changes in entropic transport costs, barycentric 
flows, and geometric curvature. The resulting framework provides a powerful and 
computationally robust tool for analyzing RDCC signatures in weak-lensing topology.
\end{abstract}
% ---------------------------------------------------------
\section{Introduction}

Persistence diagrams encode the birth and death of topological features in weak-
lensing convergence fields. Classical Wasserstein distances provide a natural 
metric on the space of diagrams, but they are computationally expensive and 
non-smooth. Entropic optimal transport introduces a regularization term that 
smooths the transport plan, yielding the Sinkhorn divergence—a fast, stable, and 
differentiable alternative.

In cosmology, these properties are particularly valuable: weak-lensing topology 
requires robust metrics that remain stable under noise, sampling variation, and 
small-scale fluctuations. In the Relaxation-Driven Cyclic Cosmology (RDCC), the 
relaxation mechanism modifies the gravitational potential in a scale-dependent 
manner, producing characteristic signatures in entropic transport geometry.
% ---------------------------------------------------------
\section{Entropic Optimal Transport}

Given two persistence diagrams $D$ and $E$ with cost matrix $C$, the entropic OT 
problem is

\begin{equation}
    \mathrm{OT}_{\varepsilon}(D,E)
    = \min_{\pi}
      \left[
        \sum_{i,j} \pi_{ij} C_{ij}
        + \varepsilon \sum_{i,j} \pi_{ij} \log \pi_{ij}
      \right],
\end{equation}

where $\varepsilon > 0$ is the entropic regularization parameter.

The Sinkhorn divergence is

\begin{equation}
    S_{\varepsilon}(D,E)
    = \mathrm{OT}_{\varepsilon}(D,E)
    - \frac{1}{2}
      \mathrm{OT}_{\varepsilon}(D,D)
    - \frac{1}{2}
      \mathrm{OT}_{\varepsilon}(E,E).
\end{equation}

This divergence is symmetric, positive definite, and smooth.
% ---------------------------------------------------------
\section{RDCC Modification of Entropic Transport}

In RDCC, the convergence evolves as

\begin{equation}
    \kappa_{\mathrm{RDCC}}(\ell)
    = \kappa(\ell)
      \left[
        1 - \chi_{\kappa}
        \left(\frac{\ell}{\ell_{\mathrm{rel}}}\right)^{\alpha}
      \right].
\end{equation}

This modifies the birth and death coordinates of topological features:

\begin{align}
    b_{\mathrm{RDCC}} &= b \left[ 1 - \chi_{b} (\ell/\ell_{\mathrm{rel}})^{\alpha} \right], \\
    d_{\mathrm{RDCC}} &= d \left[ 1 - \chi_{d} (\ell/\ell_{\mathrm{rel}})^{\alpha} \right].
\end{align}

The entropic OT cost becomes

\begin{equation}
    \mathrm{OT}^{\mathrm{RDCC}}_{\varepsilon}(D,E)
    = \mathrm{OT}_{\varepsilon}(D,E)
      \left[
        1 - \chi_{\mathrm{OT}}
        \left(\frac{\ell}{\ell_{\mathrm{rel}}}\right)^{\alpha}
      \right].
\end{equation}
% ---------------------------------------------------------
\section{Sinkhorn Geometry in RDCC}

The RDCC Sinkhorn divergence is

\begin{equation}
    S^{\mathrm{RDCC}}_{\varepsilon}(D,E)
    = S_{\varepsilon}(D,E)
      \left[
        1 - \chi_{S}
        \left(\frac{\ell}{\ell_{\mathrm{rel}}}\right)^{\alpha}
      \right].
\end{equation}

RDCC induces:

\begin{itemize}
    \item reduced entropic transport cost at small scales,
    \item enhanced cost at large scales,
    \item a characteristic turnover near $\ell_{\mathrm{rel}}$,
    \item modified Sinkhorn barycenters,
    \item altered entropic geodesics in topological space.
\end{itemize}
% ---------------------------------------------------------
\section{Applications to Cosmological Topology}

Entropic OT enables:

\begin{itemize}
    \item stable comparison of noisy persistence diagrams,
    \item fast computation of barycenters across redshift bins,
    \item differentiable topological metrics for machine learning,
    \item entropic Wasserstein flows for RDCC parameter inference,
    \item robust RDCC vs. ΛCDM discrimination under noise.
\end{itemize}

These applications provide a computationally efficient and geometrically rich 
framework for analyzing RDCC signatures in weak-lensing topology.
% ---------------------------------------------------------
\section{Conclusion}

Entropic optimal transport and Sinkhorn geometry provide a smooth, stable, and 
computationally efficient framework for analyzing persistence diagrams in the 
Relaxation-Driven Cyclic Cosmology. The relaxation mechanism modifies entropic 
transport costs and geometric curvature in a scale-dependent manner, producing 
distinctive signatures in weak-lensing topology. The present Companion establishes 
the foundation for future studies of entropic geometry, Schrödinger bridges, and 
regularized optimal transport in cosmological topology.

\bibliographystyle{apsrev4-2}
\nocite{*}
\bibliography{references}


\end{document}
